

%% ===================================================================
\section{\vrev{기술통계 종합}}
\label{sec:descriptive_summary}
%% ===================================================================

\p{[}표~\vref{tab:descriptive_summary}\p{]}는 C0--C3 4개 조건의 효율성 비율 $\eta$에
대한 기술통계량을 종합한 것이다.
\chg{여기서 C0는 표준편차가 0이어서 왜도와 초과 첨도가 정의되지 않으므로 ``---''로 표기하였다.}

\begin{table}[htbp]
  \centering
  \caption{효율성 비율 $\eta$의 조건별 기술통계 종합}
  \label{tab:descriptive_summary}
  \small
  \begin{tabular}{lccccccc}
    \toprule
    조건 & $\bar{\eta}$ & SD & $\DRTO$ (h) & 최솟값 & 최댓값 & 왜도 & 초과 첨도 \\
    \midrule
    C0 (정적) & 1.000 & 0.000 & 6.000 & 1.000 & 1.000 & --- & --- \\
    C1 (관측성) & 0.853 & 0.125 & 5.115 & 0.468 & 1.000 & $-0.84$ & $-0.20$ \\
    C2 (Gaussian) & 1.682 & 0.615 & 10.094 & 0.477 & 3.443 & $0.26$ & $-0.77$ \\
    C3 (Pareto) & 1.514 & 0.791 & 9.085 & 0.472 & 4.000 & $1.32$ & $1.03$ \\
    \bottomrule
  \end{tabular}
\end{table}

네 조건의 기술통계에서 다음의 구조적 패턴이 도출된다.

\chg{첫째, 평균의 2단계 분기 구조를 보면,} C2와 C3는 모두 C1에서 분기하는 병렬 경로이므로,
$\eta$ 평균의 변화를 공통 경로와 분기 경로로 구분하여 해석한다.
\chg{공통 경로인 C0($1.000$)에서 C1($0.853$)로는 관측성 도입에 의해 평균이 $14.7\%$ 감소하여 복구 시간이
단축된다. 분기 A인 C1($0.853$)에서 C2($1.682$)로는 가우시안 불확실성 추가 시 불확실성 프리미엄에 의해
$\erf{97.2}{97.3}\%$ 증가하며, 분기 B인 C1($0.853$)에서 C3($1.514$)로는 파레토 불확실성 추가 시 $77.5\%$ 증가하되
C2보다 낮은 평균을 산출한다.}

C2와 C3의 평균 차이($1.682$ vs.\ $1.514$)는 분포 특성의 차이를 반영한다.
가우시안 불확실성은 대칭적으로 $\eta$를 상향시키는 반면,
파레토 불확실성은 소수 극단값에 질량이 집중되어
평균은 C2보다 낮으나 분산과 극단값 범위는 더 크다.

\chg{둘째, 표준편차의 단조 증가를 보면,} SD는 공통 경로에서 C0($0$) $\to$ C1($0.125$)로 증가하며,
분기 후 C2($0.615$) 및 C3($0.791$) 모두 추가로 증가한다.
C3의 SD가 C2의 $1.29$배인 것은 파레토 분포의
극단값이 분산을 확대시킨다는 것을 정량적으로 보여준다.

\chg{셋째, 분포 형태의 질적 전환을 보면,} 왜도(skewness)의 변화가 주목할 만하다. C1은 음의 왜도($-0.84$)를
보여 $\eta$가 좌측으로 치우치며, C2는 왜도가 거의 $0$에 가까워
($0.26$) 약한 양의 비대칭을 보인다. 반면 C3는 강한 양의 왜도($1.32$)를
보여 우측 꼬리가 길게 늘어진 비대칭 분포를 형성한다.
첨도(kurtosis)도 C3에서 $1.03$으로 가장 높아, 극단값의 발생 빈도가
정규분포 대비 높음을 나타낸다.

\p{[}그림~\vref{fig:eta_distribution_all}\p{]}은 4개 조건의 $\eta$ 분포를 밀도 곡선으로
중첩하여 비교한 것이다.
\chg{그림에서 C1은 관측성 단축에 의한 좌측 편향을, C2는 평균 중심 분포(왜도 $\approx 0$)의 상향(불확실성
프리미엄)을, C3는 우측 파레토 꼬리에 의한 극단 위험 확대를 보인다. 여기서 왜도(skewness)는 분포의 비대칭
정도를 나타내어 양의 왜도는 오른쪽 꼬리가, 음의 왜도는 왼쪽 꼬리가 긴 분포를 의미하며, 초과 첨도(excess
kurtosis)는 정규분포 대비 꼬리의 무거움을 나타내어 양수는 뾰족하고 극단값이 빈번한 분포(leptokurtic),
음수는 납작하고 극단값이 드문 분포(platykurtic)를 가리킨다. C2의 입력은 가우시안이나 시그모이드 바운딩이
극단값을 $(-1,\;+1)$ 범위로 압축하여 출력 분포의 꼬리가 정규분포보다 가벼워지므로 첨도가 음수($-0.77$)로
나타나는 반면, C3는 파레토 입력의 극단값이 시그모이드 압축을 거쳐도 꼬리 특성이 잔존하여 첨도가
양수($+1.03$)를 유지한다.}

\begin{figure}[htbp]
  \centering
  \begin{tikzpicture}
    \begin{axis}[
      width=0.92\textwidth,
      height=8cm,
      xlabel={효율성 비율 $\eta$},
      ylabel={밀도 (density)},
      xmin=0.2, xmax=4.2,
      ymin=0,
      ymajorgrids=true,
      grid style={dashed, black},
      legend style={at={(0.98,0.95)}, anchor=north east, font=\small},
      legend cell align=left,
      ]
      % C0: vertical line (Dirac delta, η=1.000 고정)
      \addplot[black, ultra thick, dashed]
        coordinates {(1.0,0) (1.0,7.0)};
      \addlegendentry{C0 ($\eta = 1.000$)}

      % C1: 실데이터 KDE — 경계 보정(reflection at η=1.0)
      \addplot[blue, thick, smooth] coordinates {
        (0.425,0.005) (0.445,0.022) (0.466,0.062) (0.486,0.130) (0.506,0.219)
        (0.526,0.320) (0.546,0.425) (0.566,0.534) (0.587,0.639) (0.607,0.734)
        (0.627,0.831) (0.647,0.936) (0.667,1.049) (0.687,1.179) (0.708,1.313)
        (0.728,1.433) (0.748,1.550) (0.768,1.687) (0.788,1.857) (0.808,2.045)
        (0.829,2.221) (0.849,2.388) (0.869,2.580) (0.889,2.841) (0.909,3.196)
        (0.929,3.678) (0.950,4.377) (0.960,4.820) (0.970,5.292) (0.980,5.728)
        (0.990,6.040) (1.000,6.154)
      };
      \addlegendentry{C1 ($\bar{\eta} = 0.853$)}

      % C2: 실데이터 KDE (n=10,000, Scott bandwidth)
      \addplot[green!60!black, thick, smooth] coordinates {
        (0.320,0.004) (0.401,0.017) (0.481,0.053) (0.561,0.114) (0.641,0.194)
        (0.722,0.277) (0.802,0.356) (0.882,0.428) (0.963,0.486) (1.043,0.523)
        (1.123,0.544) (1.203,0.555) (1.284,0.553) (1.364,0.545) (1.444,0.535)
        (1.524,0.531) (1.605,0.531) (1.685,0.534) (1.765,0.526) (1.845,0.506)
        (1.926,0.486) (2.006,0.474) (2.086,0.459) (2.167,0.433) (2.247,0.406)
        (2.327,0.383) (2.407,0.359) (2.488,0.328) (2.568,0.292) (2.648,0.253)
        (2.728,0.211) (2.809,0.167) (2.889,0.127) (2.969,0.093) (3.049,0.066)
        (3.130,0.046) (3.210,0.031) (3.290,0.019) (3.371,0.011) (3.451,0.005)
      };
      \addlegendentry{C2 ($\bar{\eta} = 1.682$)}

      % C3: 실데이터 KDE (n=10,000, Scott bandwidth)
      \addplot[red, thick, smooth] coordinates {
        (0.240,0.004) (0.347,0.024) (0.454,0.089) (0.561,0.219) (0.668,0.402)
        (0.775,0.616) (0.882,0.823) (0.936,0.896) (0.989,0.932) (1.043,0.927)
        (1.096,0.885) (1.150,0.818) (1.203,0.742) (1.257,0.667) (1.310,0.598)
        (1.364,0.537) (1.417,0.484) (1.471,0.439) (1.524,0.400) (1.578,0.368)
        (1.631,0.342) (1.685,0.319) (1.738,0.298) (1.792,0.277) (1.845,0.256)
        (1.899,0.238) (1.953,0.221) (2.006,0.208) (2.060,0.197) (2.113,0.188)
        (2.167,0.179) (2.274,0.164) (2.381,0.151) (2.488,0.137) (2.595,0.121)
        (2.702,0.110) (2.809,0.104) (2.916,0.098) (3.023,0.092) (3.130,0.086)
        (3.237,0.081) (3.344,0.077) (3.451,0.074) (3.558,0.071) (3.665,0.069)
        (3.772,0.068) (3.879,0.063) (3.986,0.044) (4.093,0.019) (4.200,0.005)
      };
      \addlegendentry{C3 ($\bar{\eta} = 1.514$)}

      % eta=1 reference line
      \draw[black, thin, dashed] (axis cs:1.0,0) -- (axis cs:1.0,7);

      % 왜도·첨도 주석
      \node[blue, font=\scriptsize, align=left, anchor=west]
        at (axis cs:0.35,3.5) {왜도 $= -0.84$\\첨도 $= -0.20$};
      \draw[blue, thin, ->] (axis cs:0.52,3.3) -- (axis cs:0.62,1.0);

      \node[green!60!black, font=\scriptsize, align=left, anchor=west]
        at (axis cs:2.5,1.2) {왜도 $= 0.26$\\첨도 $= -0.77$};
      \draw[green!60!black, thin, ->] (axis cs:2.6,1.05) -- (axis cs:2.2,0.45);

      \node[red, font=\scriptsize, align=left, anchor=west]
        at (axis cs:1.6,2.0) {왜도 $= 1.32$\\첨도 $= 1.03$};
      \draw[red, thin, ->] (axis cs:1.8,1.85) -- (axis cs:1.7,0.30);
    \end{axis}
  \end{tikzpicture}
  \caption{C0/C1/C2/C3 효율성 비율 $\eta$ 분포 비교}
  \label{fig:eta_distribution_all}
\end{figure}
